\documentclass{article}
\usepackage{amsmath, amssymb, amsthm}
\usepackage{hyperref}
\usepackage{geometry}
\geometry{margin=1in}

\title{Erd\H{o}s Problem \#17}
\date{Last updated: 2025-08-31}
\author{Source: erdosproblems.com}

\begin{document}
\maketitle

\noindent\textbf{Status:} OPEN \quad \textbf{No prize}

\noindent\textbf{Tags:} number theory, primes

\medskip
\noindent\textbf{Problem Statement:}

Are there infinitely many primes $p$ such that every even number $n\leq p-3$ can be written as a difference of primes $n=q_1-q_2$ where $q_1,q_2\leq p$?

\bigskip
\noindent\textbf{Remarks:}

The first prime without this property is $97$. The sequence of such primes is A038133 in the OEIS. These are called cluster primes.

Blecksmith, Erd\H{o}s, and Selfridge \cite{BES99} proved that the number of such primes is\[\ll_A \frac{x}{(\log x)^A}\]for every $A>0$, and Elsholtz \cite{El03} improved this to\[\ll x\exp(-c(\log\log x)^2)\]for every $c<1/8$.

This is discussed in problem C1 of Guy's collection \cite{Gu04}.

\bigskip
\noindent\textbf{References:}

[BES99] Blecksmith, Richard and Erd\H{o}s, Paul and Selfridge, J. L., Cluster primes. Amer. Math. Monthly (1999), 43--48.
 
 [El03] Elsholtz, Christian, On cluster primes. Acta Arith. (2003), 281--284.
 
 [Gu04] Guy, Richard K., Unsolved problems in number theory. (2004), xviii+437.

\bigskip
\noindent\small{Source: \url{https://www.erdosproblems.com/17}}
\end{document}
